\chapter{Kết luận và Hướng phát triển}
\label{chap:future-work}

Chương này tổng kết những gì khóa luận đã đạt được, thẳng thắn chỉ ra các giới
hạn của hệ thống ở trạng thái hiện tại, rồi từ đó đề xuất các hướng phát triển
tiếp theo. Các hướng đề xuất không phải là danh sách tính năng ngẫu nhiên, mà
bám sát những \textit{ranh giới thiết kế} và \textit{giới hạn của phép đo} đã
được nêu ở Chương~\ref{chap:methodology} và Chương~\ref{chap:evaluation}: mỗi
hướng nhằm gỡ bỏ một giả định đã được cố ý đặt ra để giữ phạm vi khóa luận tập
trung vào luồng xử lý cốt lõi.

\section{Kết luận}
\label{sec:conclusion}

Khóa luận đã \textbf{tìm hiểu, thiết kế và hiện thực từ đầu} một nền tảng
blockchain có cấp phép định hướng doanh nghiệp, lấy cảm hứng từ kiến trúc
Hyperledger Fabric. Thay vì tái sử dụng một nền tảng có sẵn, chúng tôi tự cài
đặt từng thành phần lõi để làm chủ công nghệ: từ tầng mạng ngang hàng
\texttt{libp2p}, thuật toán đồng thuận Raft, máy ảo hợp đồng thông minh
WebAssembly, cho tới tầng sổ cái bất biến và world state UTXO.

Cụ thể, hệ thống hiện thực trọn vẹn triết lý tách ba trách nhiệm
\textbf{Execute--Order--Validate} qua bốn tiến trình độc lập, với các đóng góp
nổi bật: (i)~một hiện thực Raft hoàn chỉnh kèm biến thể bầu lãnh đạo theo độ ưu
tiên loại bỏ chia phiếu và hội tụ xác định; (ii)~một máy ảo WASM với cache module
và pool sandbox cho thực thi song song; (iii)~một cơ chế cắt block lai giữa kích
thước lô và thời gian, kèm giao block theo mô hình đăng-ký--phát-tỏa và đồng bộ
bốn pha dựa trên \texttt{(commitIndex, commitHash)}; và (iv)~một phương pháp đo
lường ba lớp (client, submit, commit) cho phép cô lập và phân tích trần thông
lượng của từng tầng.

Về mặt hiệu năng, thực nghiệm ở Chương~\ref{chap:evaluation} cho thấy tầng
ordering \textit{không} phải là nút thắt: trần đóng block thuần đạt
${\sim}73$--$82$k tx/s, cao hơn nhiều lần thông lượng end-to-end tham khảo qua
toàn pipeline. Kết quả này vừa xác nhận tính đúng đắn của thiết kế đồng thuận,
vừa chỉ rõ cổ chai thực tế nằm ở đường Core (WASM và ký endorsement) --- một tri
thức định hướng trực tiếp cho các hướng phát triển bên dưới.

\section{Giới hạn hiện tại}
\label{sec:limitations}

Ở trạng thái hiện tại, hệ thống còn các giới hạn sau, phần lớn là hệ quả của
những giả định được cố ý đặt ra để giữ phạm vi:

\begin{itemize}
  \item \textbf{Mô hình tin cậy dừng lỗi (crash fault tolerance).} Raft chỉ chịu
  được lỗi \textit{dừng} (crash/mạng phân mảnh), không chịu được node hành xử
  ác ý (Byzantine). Khóa luận mới chỉ \textit{thảo luận} BFT
  (Mục~\ref{sec:contributions}, Chương~\ref{chap:related-work}) chứ chưa hiện
  thực.

  \item \textbf{Ràng buộc ``một block in-flight''.} Vòng đồng thuận hiện đề xuất
  và commit tuần tự từng block một; leader phải chờ block trước được commit mới
  cắt block kế. Đây là một cận trên tiềm ẩn của thông lượng khi chạy nhiều node
  trên mạng thật (Mục~\ref{sec:eval-threats}).

  \item \textbf{Trạng thái đồng thuận in-memory, chưa bền vững.} Log Raft và
  trạng thái bầu cử hiện giữ trong bộ nhớ; chưa có ghi trước (write-ahead log),
  fsync, snapshot hay nén log. Một node khởi động lại phải dựa vào cơ chế sync
  để bắt kịp, và cả cluster không sống sót qua sự cố mất điện đồng thời.

  \item \textbf{Thành viên tĩnh.} Danh sách thành viên và độ ưu tiên được cấu
  hình tĩnh khi khởi động; chưa hỗ trợ thêm/bớt node lúc chạy (dynamic
  membership reconfiguration).

  \item \textbf{Phép đo là cận trên của một máy đơn.} Số liệu ở
  Chương~\ref{chap:evaluation} đo trên cluster một node, loopback, in-memory, và
  giao dịch giả lập không thực thi WASM. Chúng chưa phản ánh RTT đồng thuận nhiều
  node, độ trễ mạng thật và chi phí I/O đĩa.

  \item \textbf{Bảo mật danh tính và kênh truyền còn sơ khai.} Hệ thống chưa có
  hạ tầng chứng chỉ (MSP/CA) đầy đủ, kênh \texttt{libp2p} chưa siết TLS/xác thực
  hai chiều bắt buộc, và chưa có phân quyền truy cập ở mức tổ chức.
\end{itemize}

\section{Hướng phát triển}
\label{sec:future-directions}

\subsection{Nâng cấp mô hình tin cậy lên chịu lỗi Byzantine}
\label{sec:fw-bft}

Hướng có giá trị học thuật cao nhất là thay/bổ sung tầng đồng thuận CFT hiện tại
bằng một thuật toán \textbf{chịu lỗi Byzantine} (BFT). Do kiến trúc đã tách bạch
tầng ordering sau một giao diện thông điệp rõ ràng
(Mục~\ref{sec:ordering-messages} về mô hình thông điệp và định tuyến), việc thay
lõi đồng thuận không buộc phải viết lại Core hay Committing Peer. Các ứng viên:
PBFT~\cite{castro1999pbft} cổ điển ($O(N^2)$ thông điệp mỗi vòng), hoặc các biến
thể hiện đại theo hướng leader-based tuyến tính hoá chi phí. Việc này nâng mô
hình an toàn từ quorum đa số ($\lfloor N/2\rfloor + 1$) lên ngưỡng
$\lfloor (N-1)/3\rfloor$ node lỗi, phù hợp với các liên minh doanh nghiệp mà các
tổ chức thành viên không hoàn toàn tin nhau.

\subsection{Bền vững hoá và phục hồi tầng đồng thuận}
\label{sec:fw-durability}

Bổ sung \textbf{write-ahead log} có fsync cho log Raft, kèm cơ chế
\textbf{snapshot và nén log} để tránh log tăng vô hạn và rút ngắn thời gian phục
hồi. Kết hợp với cơ chế sync bốn pha hiện có (Mục~\ref{sec:deliver-sync}), một
node có thể khôi phục từ snapshot gần nhất rồi phát lại phần log còn thiếu, thay
vì đồng bộ lại từ đầu. Đây là điều kiện tiên quyết để cluster sống sót qua sự cố
mất điện đồng thời và để chạy thực tế trong môi trường sản xuất.

\subsection{Đường ống hoá việc đóng block (pipelining)}
\label{sec:fw-pipelining}

Gỡ bỏ ràng buộc ``một block in-flight'' bằng cách cho phép leader đề xuất
\textbf{nhiều block gối đầu} trước khi block trước đó được commit, tương tự cơ
chế cửa sổ trượt. Hướng này trực tiếp nhắm vào cận trên thông lượng đã nhận diện
ở Mục~\ref{sec:eval-threats} và được kỳ vọng thu hẹp khoảng cách giữa trần đóng
block đo trên một node và thông lượng thực tế nhiều node, đổi lại độ phức tạp
trong quản lý trật tự commit và xử lý khi leader đổi.

\subsection{Tối ưu đường nóng của Core}
\label{sec:fw-core}

Vì thực nghiệm cho thấy cổ chai end-to-end nằm ở Core chứ không ở ordering
(Mục~\ref{sec:eval-threats}), hướng tăng thông lượng toàn hệ thống hiệu quả nhất
là tối ưu đường Core: (i)~\textbf{song song hoá xác minh chữ ký} endorsement và
gộp lô kiểm tra; (ii)~xem xét \textbf{chữ ký tổng hợp} (ví dụ BLS) để giảm số
lần xác minh trên mỗi block; (iii)~mở rộng pool sandbox WASM và cache module theo
tải; và (iv)~đo lại toàn pipeline sau tối ưu để xác định cổ chai kế tiếp.

\subsection{Đánh giá trên cụm nhiều node và mạng thật}
\label{sec:fw-multinode}

Lặp lại bộ thực nghiệm ở Chương~\ref{chap:evaluation} trên cụm \textbf{nhiều
node phân tán} qua mạng thật, với trạng thái ghi xuống đĩa và giao dịch thực thi
WASM đầy đủ. Mục tiêu là lượng hoá phần chi phí mà các giả định hiện tại đã bỏ
qua --- RTT đồng thuận, độ trễ mạng, chi phí fsync --- và kiểm chứng giả thuyết
rằng thứ hạng tương đối giữa các cấu hình (đỉnh ở block=250--500, vùng trũng ở
block=1.000) vẫn được giữ nguyên vì chúng bắt nguồn từ cơ chế gom lô độc lập với
số node.

\subsection{Cấu hình lại thành viên lúc chạy}
\label{sec:fw-membership}

Hỗ trợ thêm/bớt node và điều chỉnh độ ưu tiên \textbf{trong lúc cluster đang
chạy} mà không cần khởi động lại, dựa trên cơ chế thay đổi cấu hình an toàn của
Raft (joint consensus). Hướng này cần bảo toàn nguyên tắc quorum tính trên tổng
số thành viên (Mục~\ref{sec:quorum}) trong suốt quá trình chuyển đổi để
tránh split-brain.

\subsection{Mở rộng theo chiều ngang: đa kênh và phân mảnh}
\label{sec:fw-scaling}

Theo mô hình \textit{channel} của Hyperledger Fabric, có thể cho phép nhiều
cluster ordering độc lập phục vụ các kênh khác nhau, mỗi kênh một sổ cái riêng.
Đây là con đường mở rộng thông lượng theo chiều ngang khi một cluster đơn đã đạt
trần, đồng thời cung cấp cách ly dữ liệu giữa các nhóm tổ chức.

\subsection{Hoàn thiện hệ sinh thái hợp đồng thông minh}
\label{sec:fw-contracts}

Bổ sung một \textbf{SDK phát triển hợp đồng} đa ngôn ngữ (mọi ngôn ngữ biên dịch
được sang WASM), cơ chế \textbf{vòng đời hợp đồng} (cài đặt, nâng cấp, thu hồi có
kiểm soát), và \textbf{đo/giới hạn tài nguyên} (gas metering) để chống hợp đồng
độc hại hoặc vòng lặp vô hạn --- tận dụng khả năng cô lập sẵn có của
\texttt{wazero}.

\subsection{Siết bảo mật và sẵn sàng vận hành}
\label{sec:fw-security}

Xây dựng hạ tầng danh tính đầy đủ (MSP/CA) cho các tổ chức thành viên, bắt buộc
TLS/xác thực hai chiều trên toàn bộ kênh \texttt{libp2p}, và bổ sung phân quyền
truy cập ở mức tổ chức. Song song, hoàn thiện quan sát hệ thống (observability)
và đóng gói triển khai (ví dụ Kubernetes) để đưa các công cụ orchestrator hiện
nằm ngoài phạm vi (Mục~\ref{sec:contributions}) lên mức sẵn sàng sản xuất.
